\documentclass[11pt]{amsart}
\usepackage{amsmath,amssymb,amsthm,enumitem}
\newtheorem{theorem}{Theorem}
\newtheorem{conjecture}[theorem]{Conjecture}
\newtheorem{proposition}[theorem]{Proposition}
\newtheorem{remark}[theorem]{Remark}
\title{$E_3$-Topological Structure for the Monster VOA\\
via the Orbifold Route}
\begin{document}
\maketitle

\section{The problem}

The Moonshine module $V^\natural$ (Frenkel--Lepowsky--Meurman) has
central charge $c = 24$, a conformal vector, and no known
gauge-theoretic origin.  Its Li-filtration associated graded is
\emph{not} freely generated ($\dim V_1^\natural = 0$), so the
Khan--Zeng theorem does not apply.

\begin{conjecture}
$Z^{\mathrm{der}}_{\mathrm{ch}}(V^\natural)$ carries an
$E_3$-topological algebra structure.
\end{conjecture}

\section{The orbifold route}

$V^\natural = V_\Lambda^+$ is the $\mathbb{Z}/2$-orbifold of the
Leech lattice VOA $V_\Lambda$.

\begin{proposition}[The invariant subalgebra is $E_3$-topological]
$V_\Lambda$ is $E_3$-topological (abelian holomorphic Chern--Simons
for $U(1)^{24}$ provides the 3d HT bulk).  The conformal vector
$T = \frac{1}{2}\sum_{i=1}^{24} {:}J^i J_i{:}$ and the antighost~$G$
are both $\mathbb{Z}/2$-invariant (quadratic in the currents).
Taking $\mathbb{Z}/2$-invariants of an $E_n$-algebra preserves the
$E_n$-structure.  Therefore $V_\Lambda^+$ inherits an
$E_3$-topological structure on its derived chiral centre.
\end{proposition}

\section{The full orbifold: twisted sectors}

$V^\natural = V_\Lambda^+ \oplus V_\Lambda^{\mathrm{tw},+}$
(untwisted + twisted sector).  The twisted sector requires
$\mathbb{Z}/2$-\emph{gauging} in the 3d bulk (summing over
$\mathbb{Z}/2$ gauge fields with twisted boundary conditions).

\begin{enumerate}
\item \textbf{3d theory.}
  The $\mathbb{Z}/2$-gauged abelian holomorphic Chern--Simons:
  $(U(1)^{24}\text{ CS})/\mathbb{Z}_2$.  The gauging is
  well-defined: the action $S = \frac{1}{2}\int A \wedge
  \bar\partial A$ is quadratic hence $\mathbb{Z}/2$-invariant,
  and the anomaly vanishes (Leech lattice is even unimodular,
  modular invariance of $V^\natural$ confirms absence of
  't~Hooft anomaly).

\item \textbf{BV quantisation.}
  The orbifold BV complex requires: $\mathbb{Z}/2$-equivariant
  BV fields, $\mathbb{Z}/2$-twisted sectors in the BV formalism,
  and verification of the BV master equation in the gauged theory.
  This has not been constructed in the Costello--Li framework but
  is expected to be tractable for abelian theories with finite
  group gauging.

\item \textbf{Boundary identification.}
  The boundary of $(U(1)^{24}\text{ CS})/\mathbb{Z}_2$ should
  be $V^\natural$ (untwisted + twisted sectors).
\end{enumerate}

\section{What this achieves}

The orbifold route reduces the conjecture for $V^\natural$ to a
bounded technical construction: orbifold BV quantisation of abelian
holomorphic Chern--Simons with finite group gauging.  The abelian
case is the simplest possible testing ground for orbifold HT theory.

\section{Beyond $V^\natural$}

The van Ekeren--M\"oller--Scheithauer theorem shows that all
strongly rational holomorphic $c = 24$ VOAs are orbifold
constructions from Niemeier lattice VOAs.  If the orbifold BV
construction works, it gives $E_3$-topological for all 71
Schellekens candidates.

The class of conformal VAs that are NOT orbifolds of
gauge-theoretic VAs may be empty in practice, though this is
not proved.

\end{document}
